% Explicit auxiliary suppliers used by arithmetic formula and matrix closure.
\subsection{Guard collection, binder input, and canonical input}
\label{sec:term-guard-auxiliaries}

The formula and induction-matrix constructions also need three auxiliary
guard suppliers.  They do not follow merely from the arbitrary-term
vector-realization theorem, so we give their algorithms explicitly.
All translations in a native guard loop use one fixed parameter tuple
$\lambda$; no recursive call changes that tuple.

Let $G_0=\top$ and $G_{j+1}=G_j\mathbin\&E_j$ be any of the
domain or realization guards.  Given an \emph{explicit} primitive
recursive array $e_j$ of proofs of $A\to T_\lambda(E_j)$, define
$c_0$ to be truth weakening and form $c_{j+1}$ by the fixed native
implication-pairing recipe from $c_j,e_j$.  The bounded invariant is
that $c_j$ is a checked proof of $A\to T_\lambda(G_j)$, conjoined
with validity of the first $j$ entries.  Quantifier-free induction on
$j$ proves the assertion.  This is the direct-list interpretation of
\texttt{NumericPAeqValueDomainGuardCollection} and
\texttt{NumericPAeqValueRealizationGuardCollection}; it does not
extract an unspecified choice function from
$\forall j<k\,\exists e_j$.

For domain prepend, use antecedent
$T_\lambda(G_k(b))\wedge T_\lambda(D(x))$ and collect $k+1$
entries.  The zeroth proof is its right projection.  Entry $j+1$
is its left projection followed by the already completed projection
$G_k(b)\to D(b_j)$.  The list equations
$(x::b)_0=x$ and $(x::b)_{j+1}=b_j$ prove the desired conclusion
$G_{k+1}(x::b)$ exactly.  This construction is
\texttt{NumericPAeqValueDomainGuardPrepend.guard\_prepend\_certificate}.
Applying it to the bound guard and retaining the free guard gives
the input-prepend supplier used by universal matrix closure.

Write $D(x)=\operatorname{domainCode}(x)$ and
$V_x=\operatorname{variableValue}(x,x)$.
The requested self-variable proof is
$T_\lambda([D(x)\Rightarrow V_x])$.
Its fixed relational ND template proves
$D(x)\to(D(x)\wedge\operatorname{Eq}(x,x))$ from the
closed relational reflexivity sentence.  Thus it is precisely an
instance of the lower singleton-requested-template constructor,
with its sole lower nonlogical leaf supplied by the closed
\texttt{reflexivityNativeProof}; it needs no term recursion.
This is the exact boundary in
\texttt{NumericPAeqValueSelfVariableCertificate.lean}.

For binder input, let
$G=\operatorname{vectorInput}(s,r,b,f,a,w)$ and
\[
 A=G\mathbin\&D(x),\qquad
 B=\operatorname{vectorInput}
       (s+1,r+1,x::b,f,x::a,\operatorname{qVec}(w)),
 \quad S=[D(x)\Rightarrow V_x].
\]
The exact requested-label map in
\texttt{NumericPAeqSupportedVectorBinderInputCertificate.lean} is
\[
 d=\max(S,[A\Rightarrow B]),\qquad
 \lambda'=P_d(\lambda).
 \tag{T15}
\]
The only non-native request is the self-variable proof at $\lambda'$.
At this same tuple the new realization guard has entry zero $V_x$
and entry $j+1$ equal to the old entry $j$.  These are
\texttt{vector\_qVec\_entry\_zero} and
\texttt{vector\_qVec\_entry\_succ}; the successor equation is
the completed source-lift/value-syntax substitution identity, with
the exact support conditions of the vector environment.
Collect the zeroth entry using right projection and the self proof;
collect each later entry using left projection and old-guard
projection.  Prepend the domain guards of $b$ and $a$ by the
preceding native loop, retain the two copies of the free-domain
guard, and pair the resulting five factors in the literal
\texttt{vectorInput} association.  Native/source implication
normalization and padding readback give $T_\lambda([A\Rightarrow B])$.
All entry indices are bounded by $s+1$ or $r+1$ and every proof
value is computed explicitly.  When formula realization first
shifts $b,f,a$, this same construction is applied to those shifted
vectors; the support and validity tests are preserved by the
already completed target-shift syntax equations.

For canonical input, let $|B|=m$ and set
\[
 G=G_m(B),\quad
 J=\operatorname{vectorInput}(m,0,0,B,B,\operatorname{fvarVec}(m)),
 \qquad S_i=[D(B_i)\Rightarrow\operatorname{variableValue}(B_i,B_i)].
\]
Replace the source's collection-based existential bound by the
explicit primitive recursion
\[
 d_0=0,\qquad d_{i+1}=\max(d_i,S_i),\qquad
 d=d_m,\qquad \lambda'=P_d(\lambda).
 \tag{T16}
\]
The bounded invariant $\bigwedge_{i<j}S_i\le d_j$ supplies every
coverage check, including the empty case $m=0$.
At $\lambda'$ request the self-variable proof for each $i<m$;
compose it with the native projection $G\to D(B_i)$.
The canonical realization entry is literally
$\operatorname{variableValue}(B_i,B_i)$, because
$\operatorname{fvarVec}(m)_i=\operatorname{fvar}(i)$ and the
value-syntax free-variable equation selects $B_i$.
Collect these $m$ explicit arrows with the first loop.
The literal input shape is
\[
 T_{\lambda'}(J)=
 (\top\wedge T_{\lambda'}G)\wedge
   ((T_{\lambda'}G\wedge T_{\lambda'}G)\wedge
       T_{\lambda'}G_{\rm realized}).
\]
Truth weakening, three identities on $G$, and the collected arrow
therefore prove $G\to J$ by fixed conjunction operations.
Native/source normalization and padding readback finish the
requested proof.  This is the explicit realization of
\texttt{NumericPAeqCanonicalVectorInputCertificate.lean} and its
\texttt{self\_codes\_bounded} auxiliary.  It neither invokes an
unbounded collection principle nor recurses through an arbitrary
formula-realization supplier.

Consequently these auxiliaries belong immediately after guard
projection and the lower fixed self-variable template, before
arithmetic formula realization and induction-matrix closure.
Their only new closed leaf is ordinary relational reflexivity,
already present among the six fixed ordinal leaves of the generic
relational construction.  No edge returns to the arithmetic
formula or proof strata.
